\documentclass[12pt,conference,onecolumn]{IEEEtran}

\usepackage[hidelinks]{hyperref}

\title{The Bridges of Monmouth County: Internship at Monmouth County Engineering}
\author{%
\IEEEauthorblockN{Ryan Leung}\IEEEauthorblockA{Science \& Engineering\\Manalapan High School\\Englishtown, NJ\\\href{mailto:426rleung@frhsd.com}{426rleung@frhsd.com}}}
\date{June 16, 2026}

\newcommand{\keywords}{Monmouth County Engineering, civil engineering, internship, bridges}

\usepackage{hyperref}
\makeatletter
\AtBeginDocument{
\hypersetup{%
pdftitle={\@title},
pdfauthor={Ryan Leung},
pdfkeywords={\keywords}}}
\makeatother

\begin{document}
\maketitle 

\begin{abstract}
Monmouth County Engineering’s Bridge Department is responsible for the construction, maintenance, and improvement of county-owned bridges. Of the aforementioned duties, the county’s main duty is to maintain its already existing bridges and ensure the maintenance of each bridge goes smoothly. As most of the design-work and calculations are done out-of-house, the county is mainly responsible for the allocation and proper compensation of the work to private contractors.

During my internship at the county, I gained hands-on experience of each step of a bridge project that most private engineering firms would be unable to get, as the county oversees the entire process. I was able to learn about reading and interpreting construction plans, the different techniques used to ease construction, the use of AutoCAD, and the correspondence required to make these jobs possible. All of this was aided by my mentors’ guidance and frequent visits to their many job-sites. Some of these job-sites include the bridges U-33, MT-10, A-38, S-15, and S-13.
\end{abstract}

\begin{IEEEkeywords}
\keywords
\end{IEEEkeywords}

\end{document}
